%===== General Introduction =====%

\chapter*{General Introduction}
\addcontentsline{toc}{chapter}{General Introduction}
\markboth{GENERAL INTRODUCTION}{}

\section*{Context}

The rapid growth of internet access and smartphone adoption in Tunisia has significantly transformed consumer purchasing habits and accelerated the development of online commerce. More customers now rely on digital platforms to search for products, compare prices, and place orders remotely. In this evolving environment, companies specialized in technology products increasingly require modern digital solutions to remain competitive and improve customer experience.

\textbf{Malltech} is a Tunisian company specialized in selling technological products such as smartphones, electronics, accessories, and various tech equipment. Despite the growing demand for online shopping, the company lacked a complete digital platform capable of centralizing product management, customer orders, inventory monitoring, and online sales operations.

This project was carried out during an internship at \textbf{Anter Lab} with the objective of designing and developing a complete e-commerce solution for Malltech. The platform includes both a customer-facing e-commerce website and a comprehensive administration backoffice to support and optimize the company’s commercial activities.

\section*{Problem Statement}

Many local businesses still depend on fragmented management methods such as social media pages, manual order processing, spreadsheets, or limited management systems. These approaches often lead to difficulties in inventory tracking, order management, customer follow-up, and business scalability.

Existing international e-commerce solutions may provide generic functionalities, but they often require expensive customization and are not always adapted to the Tunisian market and local business needs. In addition, some solutions lack flexibility, centralized management, or advanced administrative features.

The main challenge of this project is therefore to develop a modern, scalable, and centralized e-commerce platform that allows Malltech to:
\begin{enumerate}
    \item Present and sell products online through a modern and responsive storefront.
    \item Centralize product, inventory, order, and customer management.
    \item Simplify administrative and operational tasks through a complete backoffice.
    \item Improve customer experience and strengthen the company’s digital presence.
    \item Provide a scalable architecture that can evolve with future business requirements.
\end{enumerate}

\newpage

\section*{Project Objectives}

\textbf{Primary objectives:}
\begin{itemize}
    \item Design and develop a complete e-commerce website for online product sales.
    \item Build a comprehensive administration backoffice for managing products, categories, inventory, customers, and orders.
    \item Implement a secure authentication and authorization system for customers and administrators.
    \item Provide an optimized and responsive user experience across different devices.
\end{itemize}

\textbf{Secondary objectives:}
\begin{itemize}
    \item Integrate Search Engine Optimization (SEO) features to improve online visibility.
    \item Implement analytics and reporting functionalities for business monitoring.
    \item Develop dynamic content management capabilities for banners, pages, and promotions.
    \item Automate deployment and development workflows using modern DevOps tools.
    \item Ensure platform scalability, maintainability, and security.
\end{itemize}

\section*{Scope and Limitations}

\textbf{In scope:} Development of the e-commerce website, administration backoffice, product and inventory management modules, order management system, customer management, SEO features, analytics tools, and deployment infrastructure.

\textbf{Out of scope:} Native mobile applications, multi-vendor marketplace functionality, advanced AI recommendation systems, and external payment gateway integrations not required during the initial project phase.